\documentclass[UTF8,11pt]{ctexart}

\usepackage[a4paper,margin=2.4cm]{geometry}
\usepackage{booktabs}
\usepackage{longtable}
\usepackage{array}
\usepackage{amsmath}
\usepackage{xcolor}
\usepackage{hyperref}
\usepackage{listings}

\hypersetup{
  colorlinks=true,
  linkcolor=blue!60!black,
  urlcolor=blue!60!black
}

\lstdefinestyle{codestyle}{
  basicstyle=\ttfamily\small,
  keywordstyle=\color{blue!70!black},
  commentstyle=\color{green!40!black},
  stringstyle=\color{red!60!black},
  breaklines=true,
  columns=fullflexible,
  frame=single,
  framerule=0.3pt,
  rulecolor=\color{black!30},
  showstringspaces=false,
  tabsize=2
}
\lstset{style=codestyle}

\title{CS336 Assignment 2 Systems 实现总结}
\author{李东泽}
\date{\today}

\begin{document}
\maketitle

\section{完成内容概览}

本次完成了 Assignment 2 中公开测试覆盖的全部系统组件，并将测试适配器连接到
\path{cs336_systems} 模块。主要新增或修改文件如下：

\begin{longtable}{>{\raggedright\arraybackslash}p{0.36\linewidth}
                  >{\raggedright\arraybackslash}p{0.56\linewidth}}
\toprule
文件 & 作用 \\
\midrule
\path{cs336_systems/flash_attention.py}
& 实现 PyTorch autograd 版本 FlashAttention，以及调用 Triton kernel 的 CUDA forward/backward 版本。\\
\path{cs336_systems/distributed.py}
& 实现 individual-parameter DDP 和 bucketed DDP，包括初始化广播、反向 hook、梯度同步和 bucket scatter。\\
\path{cs336_systems/sharded_optimizer.py}
& 实现 optimizer state sharding：每个 rank 只维护一部分参数的 optimizer state，更新后广播参数。\\
\path{tests/adapters.py}
& 将官方测试入口改为薄适配层，返回或调用 \texttt{cs336\_systems} 中的实现。\\
\path{cs336_systems/__init__.py}
& 增加源码目录直接运行时的版本号 fallback，避免未安装包时导入失败。\\
\bottomrule
\end{longtable}

\section{FlashAttention}

\subsection{原理}

标准 attention 先显式构造分数矩阵
\[
S = \frac{QK^\top}{\sqrt{d}},
\quad
P = \operatorname{softmax}(S),
\quad
O = PV.
\]
显式保存完整的 $S$ 和 $P$ 会占用 $O(BN^2)$ 显存。FlashAttention 的核心思想是用分块方式计算
softmax，并只保存反向传播所需的 log-sum-exp：
\[
L_i = \log \sum_j \exp(S_{ij}).
\]
在反向传播时可以通过
\[
P_{ij} = \exp(S_{ij} - L_i)
\]
重构 softmax 概率，不需要保存完整概率矩阵。反向传播中使用：
\[
dV = P^\top dO,
\quad
dS = P \odot \left(dP - \sum_j P_{ij} dP_{ij}\right),
\quad
dQ = dS K / \sqrt{d},
\quad
dK = dS^\top Q / \sqrt{d}.
\]

本实现包含两个入口：

\begin{itemize}
  \item \texttt{FlashAttentionPytorch}：只使用 PyTorch tensor op，实现清晰、便于验证。
  \item \texttt{FlashAttentionTriton}：在 CUDA tensor 上调用自定义 Triton forward/backward kernel；无 Triton 或 CPU 环境下 fallback 到 PyTorch 版本。
\end{itemize}

\subsection{关键代码}

PyTorch forward 保存 \texttt{q,k,v,lse}，其中测试要求恰好保存一个形状为
\texttt{(batch, n\_queries)} 的 log-sum-exp tensor：

\begin{lstlisting}[language=Python]
scores = _attention_scores(q, k, bool(is_causal))
lse = torch.logsumexp(scores, dim=-1)
p = torch.exp(scores - lse.unsqueeze(-1))
out = torch.matmul(p, v)

ctx.save_for_backward(q, k, v, lse)
\end{lstlisting}

Triton forward kernel 以每个 \texttt{batch/query} 为一个 program，计算一行 attention：

\begin{lstlisting}[language=Python]
scores = tl.sum(k * q[None, :], axis=1) * scale
scores = tl.where(valid, scores, -1.0e6)

row_max = tl.max(scores, axis=0)
exp_scores = tl.exp(scores - row_max)
exp_sum = tl.sum(exp_scores, axis=0)
lse = tl.log(exp_sum) + row_max
p = exp_scores / exp_sum

out = tl.sum(p[:, None] * v, axis=0)
\end{lstlisting}

Triton backward kernel 中，\texttt{dq} 属于当前 query，可以直接写回；\texttt{dk/dv} 被多个 query 共同贡献，因此使用 atomic add：

\begin{lstlisting}[language=Python]
dp = tl.sum(v * do[None, :], axis=1)
delta = tl.sum(p * dp, axis=0)
ds = p * (dp - delta)

dq = tl.sum(ds[:, None] * k, axis=0) * scale
dk = ds[:, None] * q[None, :] * scale
dv = p[:, None] * do[None, :]

tl.atomic_add(grad_k_ptr + ..., dk, sem="relaxed", mask=...)
tl.atomic_add(grad_v_ptr + ..., dv, sem="relaxed", mask=...)
\end{lstlisting}

\section{Individual-Parameter DDP}

\subsection{原理}

Distributed Data Parallel 的目标是让每个 rank 处理不同 mini-batch shard，但参数更新等价于单机大 batch。
具体做法是：

\begin{enumerate}
  \item 初始化时从 rank 0 广播完整模型状态，确保所有 rank 参数一致。
  \item 每个可训练参数注册 backward hook。
  \item 某个参数的梯度就绪后立即执行 \texttt{all\_reduce(sum)}。
  \item optimizer step 前等待所有通信完成，并将梯度除以 \texttt{world\_size}。
\end{enumerate}

这样每个 rank 的梯度变为全局平均梯度，因此每个 rank 的 optimizer step 与单机全量 batch 一致。

\subsection{关键代码}

\begin{lstlisting}[language=Python]
for tensor in self.module.state_dict().values():
    _broadcast_(tensor, src=0)

for param in _unique_trainable_parameters(self.module):
    param.register_post_accumulate_grad_hook(self._make_hook(param))
\end{lstlisting}

\begin{lstlisting}[language=Python]
def hook(_param):
    if param.grad is None or self._world_size == 1:
        return
    work = _all_reduce_sum_(param.grad, async_op=True)
    self._handles.append((work, param.grad))

def finish_gradient_synchronization(self):
    for work, grad in self._handles:
        work.wait()
        grad.div_(self._world_size)
    self._handles.clear()
\end{lstlisting}

实现中还处理了 \texttt{gloo + CUDA tensor} 的情况：官方测试使用 gloo backend，但 CUDA 可见时模型会被放到 GPU。
对于 gloo 不直接支持的 CUDA tensor，代码先复制到 CPU 做 collective，再复制回原设备。

\section{Bucketed DDP}

\subsection{原理}

Individual-parameter DDP 每个参数一次通信，通信次数多。Bucketed DDP 将多个参数梯度拼接成 bucket，
在 bucket 内所有参数梯度都就绪后发起一次 all-reduce，从而减少通信启动开销。

本实现步骤如下：

\begin{enumerate}
  \item 根据 \texttt{bucket\_size\_mb} 按参数顺序建立 bucket。
  \item 每个参数的 hook 标记所在 bucket 中该参数已经 ready。
  \item 当 bucket 中所有参数 ready 后，将梯度 flatten 后拼接为一个 tensor，并异步 all-reduce。
  \item optimizer step 前等待通信完成，将平均后的 flat gradient scatter 回各参数的 \texttt{.grad}。
\end{enumerate}

\subsection{关键代码}

构建 bucket：

\begin{lstlisting}[language=Python]
for param in params:
    param_size = param.numel() * param.element_size()
    if current and max_size is not None and current_size + param_size > max_size:
        self._append_bucket(current)
        current = []
        current_size = 0
    current.append(param)
    current_size += param_size
\end{lstlisting}

bucket ready 后发起同步：

\begin{lstlisting}[language=Python]
bucket.ready.add(id(param))
if len(bucket.ready) == len(bucket.params) and bucket.work is None:
    self._launch_bucket(bucket)

grads = [param.grad.reshape(-1) for param in bucket.params]
bucket.flat_grad = torch.cat(grads)
bucket.work = _all_reduce_sum_(bucket.flat_grad, async_op=True)
\end{lstlisting}

同步完成后 scatter 回原参数：

\begin{lstlisting}[language=Python]
bucket.work.wait()
bucket.flat_grad.div_(self._world_size)
offset = 0
for param in bucket.params:
    numel = param.numel()
    param.grad.copy_(bucket.flat_grad[offset : offset + numel].view_as(param))
    offset += numel
\end{lstlisting}

\section{Sharded Optimizer}

\subsection{原理}

普通 AdamW 在每个 rank 上都保存所有参数的 optimizer state，例如一阶动量和二阶动量。
Optimizer state sharding 的目标是让每个 rank 只保存一部分参数的 optimizer state，降低每张卡的内存压力。

本实现采用简单稳定的参数分片策略：
\[
\operatorname{owner}(i) = i \bmod \operatorname{world\_size}.
\]
每个 rank 只把自己负责的参数传给本地 optimizer。执行 \texttt{step()} 后，每个参数由其 owner rank 广播到所有 rank，
从而恢复所有模型副本的一致性。

\subsection{关键代码}

\begin{lstlisting}[language=Python]
def _owner(self, param_index: int) -> int:
    return param_index % self.world_size

local_group["params"] = [
    param for param in group["params"]
    if self._owner(self.param_to_index[id(param)]) == self.rank
]
self.local_optimizer = optimizer_cls(local_groups, **kwargs)
\end{lstlisting}

\begin{lstlisting}[language=Python]
if self.local_optimizer is not None:
    self.local_optimizer.step()

for index, param in enumerate(self.global_params):
    _broadcast_(param.data, src=self._owner(index))
\end{lstlisting}

\section{测试数据与结果}

\subsection{测试环境}

最终 GPU 测试没有使用新建的 \texttt{uv} 环境，因为服务器根分区空间不足，安装 CUDA 版 \texttt{torch} 时失败。
改用服务器已有的 \texttt{rk-gemm} conda 环境，其 PyTorch 版本与作业依赖匹配：

\begin{itemize}
  \item Python 环境：\texttt{/home/ldz/miniconda3/envs/rk-gemm/bin/python}
  \item PyTorch：\texttt{2.6.0+cu124}
  \item GPU 限制：\texttt{CUDA\_VISIBLE\_DEVICES=0,2,4,5}
  \item 可见 GPU：4 张 \texttt{NVIDIA GeForce RTX 3090}
\end{itemize}

\subsection{公开测试结果}

\begin{longtable}{>{\raggedright\arraybackslash}p{0.20\linewidth}
                  >{\raggedright\arraybackslash}p{0.58\linewidth}
                  >{\raggedright\arraybackslash}p{0.14\linewidth}}
\toprule
测试 & 覆盖内容 & 结果 \\
\midrule
\texttt{CUDA full tests}
& FlashAttention PyTorch/Triton、bucketed DDP、individual DDP、sharded optimizer，全 CUDA 可见路径。 & 16 passed \\
\texttt{CPU-only tests}
& CPU-only 路径；CUDA/Triton attention 测试按 pytest 条件跳过。 & 12 passed, 4 skipped \\
\texttt{compileall}
& Python 语法和 import 基本检查。 & passed \\
\bottomrule
\end{longtable}

完整测试命令如下：

\begin{lstlisting}
CUDA_VISIBLE_DEVICES=0,2,4,5 \
  /home/ldz/miniconda3/envs/rk-gemm/bin/python -m pytest tests -q

CUDA_VISIBLE_DEVICES= python -m pytest tests -q

python -m compileall cs336_systems tests/adapters.py
\end{lstlisting}

CUDA 全套测试输出摘要：

\begin{lstlisting}
tests/test_attention.py::test_flash_forward_pass_pytorch PASSED
tests/test_attention.py::test_flash_forward_pass_triton[False] PASSED
tests/test_attention.py::test_flash_forward_pass_triton[True] PASSED
tests/test_attention.py::test_flash_backward_pytorch PASSED
tests/test_attention.py::test_flash_backward_triton[False] PASSED
tests/test_attention.py::test_flash_backward_triton[True] PASSED
tests/test_ddp.py::* PASSED
tests/test_ddp_individual_parameters.py::* PASSED
tests/test_sharded_optimizer.py::* PASSED

16 passed, 1 warning in 47.17s
\end{lstlisting}

\subsection{Benchmark 数据记录}

仓库中已有当前硬件的 RTX 4090 benchmark 结果，位于
\path{benchmark_outputs/2026-04-27_rtx4090_formal/RESULTS.md}。
其中 FP32 baseline timing 如下，可作为系统优化报告中的性能参考数据：

\begin{longtable}{lrrrr}
\toprule
模型 & fwd均值/ms & fwd标准差/ms & fwd+bwd均值/ms & fwd+bwd标准差/ms \\
\midrule
small & 11.448 & 0.066 & 35.383 & 0.514 \\
medium & 23.370 & 0.088 & 68.459 & 1.153 \\
large & 39.026 & 0.123 & 120.142 & 8.285 \\
xl & 73.647 & 0.129 & 200.166 & 1.484 \\
2.7b & 107.876 & 1.129 & OOM & -- \\
\bottomrule
\end{longtable}

这组数据说明：在 24GB RTX 4090 上，2.7B 模型的 forward 可以运行，但 forward+backward 在 batch size 4、
context length 128、vocab size 10000 的设定下会 OOM。

\section{结论}

Assignment 2 的核心系统组件已经完成并通过公开测试。实现上遵循了作业要求：

\begin{itemize}
  \item FlashAttention 保存 LSE 并在 backward 中重构 softmax；CUDA 路径实际调用 Triton kernel。
  \item DDP 初始化广播模型参数，并在 backward 阶段同步梯度，保证与单机大 batch 训练一致。
  \item Bucketed DDP 将多个参数梯度合并通信，降低 per-parameter all-reduce 的启动开销。
  \item Sharded optimizer 将 optimizer state 分散到不同 rank，step 后广播参数维持副本一致性。
\end{itemize}

最终验证结果为 CUDA 可见环境下 \texttt{16 passed}，CPU-only 环境下 \texttt{12 passed, 4 skipped}。

\end{document}
